\section{\textsc{SciExplore}}
% To operationalize the desiderata above, \textsc{SciExplore} instantiates scientific information-seeking as four complementary tasks. We discuss the task definitions in Section~\ref{sec:task_definition}, data construction in Section~\ref{sec:dataset_construction}, and quality control in Section~\ref{sec:quality_control}. Detailed data statistics and evaluation metrics can be found in Appendix~\ref{appendix_dataset_distribution} and Appendix~\ref{app:evaluation_details}, respectively.

To operationalize the desiderata above, \textsc{SciExplore} instantiates scientific information-seeking as four complementary tasks, curated through a fully expert-driven construction and verification pipeline. We formalize this task taxonomy in Section~\ref{sec:task_definition}, detail the corresponding data construction procedure in Section~\ref{sec:dataset_construction}, and describe our quality control mechanisms in Section~\ref{sec:quality_control}. Detailed data statistics and evaluation metrics are further provided in Appendix~\ref{appendix_dataset_distribution} and Appendix~\ref{app:evaluation_details}, respectively.

% To operationalize the desiderata above, \textsc{SciExplore} instantiates scientific information-seeking as four complementary tasks---spanning database navigation, ambiguous literature retrieval, missing reference completion, and cross-source structured knowledge synthesis---that jointly stress an agent's ability to locate, disambiguate, and integrate evidence from authentic scientific sources. Every instance is authored and verified by domain experts, and the full pipeline is designed to enforce realism, answer uniqueness, and resistance to parametric shortcuts, so that strong performance necessarily reflects genuine retrieval and reasoning rather than memorization. 
% We discuss the task definitions in Section~\ref{sec:task_definition}, data construction in Section~\ref{sec:dataset_construction}, and quality control in Section~\ref{sec:quality_control}. Detailed data statistics and evaluation metrics can be found in Appendix~\ref{appendix_dataset_distribution} and Appendix~\ref{app:evaluation_details}, respectively.

% \subsection{Task Definition and Construction}

\subsection{Task Definition}\label{sec:task_definition}


% Motivated by the limitations of existing benchmarks discussed above, \textsc{SciExplore} targets scientific information-seeking behaviors that arise in authentic research workflows but are largely underexplored in prior evaluations, and frames scientific inquiry as a multi-stage process involving database navigation, literature discovery under ambiguity, evidence consolidation, and structured synthesis.

% To operationalize these capabilities, we decompose scientific information seeking into four representative task types, each capturing a distinct stage of the scientific cognitive process with increasing complexity.
% The four task types are defined as follows.

\begin{comment}
\textsc{SciExplore} conceptualizes scientific information seeking as a progressive reasoning process that unfolds across multiple levels of abstraction.
Starting from the resolution of individual entities, we gradually increase the evaluation difficulty from document-level identification and evidence-level grounding to domain-level knowledge synthesis, and instantiate each capability through a corresponding task formulation.

Specifically, entity-level reasoning is examined through scientific database navigation, where agents must explore and traverse structured resources in the natural sciences.
Document-level identification and evidence-level grounding are evaluated via ambiguous literature retrieval and missing reference completion tasks, which require models to identify relevant papers under uncertainty and correctly align scientific claims with their supporting evidence.
Finally, domain-level synthesis is assessed through cross-source structured knowledge synthesis, where agents are required to integrate information from multiple heterogeneous studies and organize the resulting knowledge into a structured representation.
\end{comment}

\begin{comment}
As illustrated in Table \ref{tab:sciexplore-task-taxonomy}, this benchmark for evaluating science information seeking encompasses four progressively complex task types, each corresponding to different capability levels and atomic capabilities essential for effective information retrieval and synthesis in scientific research.

The first two tasks in our benchmark focus on database navigation and retrieval within scientific literature, respectively. Specifically, the Scientific Database Navigation task involves entity-level reasoning, assisting researchers to identify, filter, and disambiguate relevant entities while navigating through complex databases. This is crucial for researchers who must efficiently locate specific data amidst vast resources. Complementing this, the Ambiguous Literature Retrieval task targets document-level identification by helping users disambiguate intent and effectively discriminate among various documents to eliminate irrelevant ones. These tasks work together to evaluate the ability of LLMs or agents to find precise information in scientifically dense environments.

The third task, Missing Reference Completion, emphasizes evidence-level grounding. This task is essential for ensuring the integrity of claims made in scientific discourse. It involves processes such as claim abstraction, aligning claims with supporting evidence, and discovering and validating relevant evidence. This task not only aids in completing references but also enriches the overall quality of scientific research by ensuring that claims are adequately supported.

Finally, the Cross-source Structured Knowledge Synthesis task operates at the domain-level synthesis, representing a higher-order capability in our benchmark. This task involves domain abstraction and the integration of information from multiple sources to facilitate structured synthesis. It is critical for scholars who aim to produce comprehensive reviews or syntheses of existing literature, as it allows for the coherent amalgamation of knowledge across various studies, thus fostering a holistic understanding of specific scientific domains.
\end{comment}
% Table~\ref{tab:sciexplore-task-taxonomy} elaborates elaborates on the taxonomy of tasks in \textsc{SciExplore}.
% % As illustrated in Table~\ref{tab:sciexplore-task-taxonomy}, this benchmark for evaluating science information seeking includes four progressively complex task types, each corresponding to varying abilities essential for effective information retrieval and synthesis in scientific research.
% Specifically, T1, Scientific Database Navigation, helps researchers identify, filter, and disambiguate relevant entities in complex databases, which is crucial for efficiently locating specific data. T2, Ambiguous Literature Retrieval, aids users in disambiguating intent and discriminating among documents to eliminate irrelevant ones. Together, these tasks assess the ability of LLMs or agents to locate precise information in dense scientific environments.

% T3, Missing Reference Completion, ensures the integrity of scientific claims. This task includes claim abstraction, aligning claims with supporting evidence, and validating relevant evidence, thereby enriching the quality of research by ensuring adequate support for claims. Finally, T4, Cross-source Structured Knowledge Synthesis, represents a higher-order capability, involving domain abstraction, information extraction and the integration of information from multiple sources for structured synthesis. 


Table~\ref{tab:sciexplore-task-taxonomy} elaborates on the taxonomy of tasks in \textsc{SciExplore}, which spans a progression from entity-level retrieval to cross-source structured synthesis and is designed to jointly cover the core information-seeking competencies required in authentic scientific research workflows. 

Specifically, T1, Scientific Database Navigation, helps researchers identify, filter, and disambiguate relevant entities in complex databases, which is crucial for efficiently locating specific data. T2, Ambiguous Literature Retrieval, aids users in disambiguating intent and discriminating among documents to eliminate irrelevant ones. Together, these tasks assess the ability of LLMs or agents to locate precise information in dense scientific environments.
T3, Missing Reference Completion, ensures the integrity of scientific claims. This task includes claim abstraction, aligning claims with supporting evidence, and validating relevant evidence, thereby enriching the quality of research by ensuring adequate support for claims. Finally, T4, Cross-source Structured Knowledge Synthesis, represents a higher-order capability, involving domain abstraction, information extraction, and the integration of information from multiple sources for structured synthesis. Collectively, T3 and T4 probe the evidence-grounding and cross-source integration competencies that go beyond mere information location, complementing T1--T2 and together forming a comprehensive evaluation suite that mirrors the end-to-end workflow of a scientific research assistant.

% This task is critical for scholars aiming to produce comprehensive reviews or syntheses, as it facilitates the coherent amalgamation of knowledge from various studies, fostering a holistic understanding of scientific domains.


% Scientific information seeking in real research workflows is inherently progressive, requiring agents to advance from structured entity retrieval to document discovery, contextual evidence grounding, and ultimately domain-level knowledge synthesis.
% Guided by this observation, \textsc{SciExplore} is designed as a unified evaluation framework in which task definitions and data construction are jointly specified, so that each task instance instantiates a distinct level of scientific reasoning difficulty.
% All tasks are manually constructed by domain experts to ensure realism, answer uniqueness, and resistance to shortcut retrieval.

% \textsc{SciExplore} comprises four task types, ordered by increasing abstraction and integration demands, corresponding to successive stages of scientific information seeking. Details of the full dataset construction process are provided in the Appendix~\ref{app:datasets_construction}.

% \noindent{\textbf{Scientific Database Navigation.}}
% At the entity level, agents must recover uniquely identifiable targets through structured reasoning over interconnected scientific databases.
% To enforce genuine multi-hop traversal, task instances are constructed using a \emph{Reverse Trajectory Construction Strategy}, in which explicit entities are progressively masked and replaced with nested attribute-based constraints, such that correct resolution requires reversing the underlying dependency chain.

% \noindent{\textbf{Ambiguous Literature Retrieval.}}
% At the document level, agents are required to identify a specific target paper from imprecise and noisy descriptions, reflecting exploratory literature search under uncertainty.
% Construction adopts a two-stage design that combines \emph{Feature Denoising and Fuzzification} to substantially enlarge the retrieval space with \emph{Validation Constraint Injection} to preserve answer uniqueness through post-retrieval verification cues.

% \noindent{\textbf{Missing Citation Completion.}}
% At the argument level, evaluation shifts from document identification to contextual evidence grounding.
% Given a scientific paragraph with missing citations, agents must infer which prior works substantiate individual claims.
% To prevent shortcut matching, task contexts are rewritten to minimize surface-level overlap with the target references, ensuring that successful citation reconstruction depends on reasoning over claim--evidence relationships.

% \noindent{\textbf{Cross-source Structured Knowledge Synthesis.}}
% At the domain level, agents are challenged to integrate evidence across multiple studies into a structured representation.
% Given a research topic and a predefined comparison schema, tasks are constructed from expert-curated sources with strict format and schema constraints, such that accurate completion requires precise cross-paper evidence integration rather than partial or approximate extraction.

\begin{table*}[t]
\centering
\caption{Task taxonomy of \textsc{SciExplore}, detailing the task types, corresponding capability levels, atomic capabilities, task counts, and associated metrics.}
\label{tab:sciexplore-task-taxonomy}
% \vspace{-0.5em}
\scriptsize
% \footnotesize
\begin{tabular*}{\textwidth}{l @{\extracolsep{\fill}} l @{} l @{} l @{} c @{} l}
\toprule
\textbf{No.} & \textbf{Task Type} & \textbf{Capability Level} & \textbf{Atomic Capabilities} & \textbf{Tasks} & \textbf{Metric} \\
\midrule
T1 & \makecell[l]{Scientific Database \\ Navigation} & \makecell[l]{Entity-Level \\Reasoning} & \makecell[l]{Entity identification, filtering, disambiguation, \\ multi-constraint navigation} & 39 & Accuracy \\
\midrule
T2 & \makecell[l]{Ambiguous Literature \\ Retrieval} & \makecell[l]{Document-Level\\Identification} & \makecell[l]{Intent disambiguation, \\ document-level discrimination and elimination} & 32 & Accuracy \\
\midrule
T3 & \makecell[l]{Missing Reference \\Completion} & \makecell[l]{Evidence-Level\\Grounding} & \makecell[l]{Claim abstraction, claim--evidence alignment, \\ evidence discovery and validation} & 14 & Accuracy \\
\midrule
T4 & \makecell[l]{Cross-source Structured \\ Knowledge  Synthesis} & \makecell[l]{Domain-Level \\ Synthesis} & \makecell[l]{Domain abstraction, information extraction, \\ cross-source integration, structured synthesis} & 18 & Recall \\
\bottomrule
\end{tabular*}
% \vspace{-1em}
\end{table*}

\subsection{Dataset Construction}
\label{sec:dataset_construction}

\textsc{SciExplore} is constructed through a fully expert-driven process to ensure realism, answer uniqueness, and resistance to shortcut retrieval.
All task instances are manually curated by domain experts and designed to reflect authentic scientific research behaviors, rather than synthetic or template-based generation, as shown in Fig.~\ref{fig:dataset-pipeline}.

\vspace{0.5em}
\noindent{\textbf{Scientific Database Navigation.}} For T1, we design tasks that require genuine multi-step database navigation using a \emph{Reverse Trajectory Construction Strategy}, where explicit entity identifiers are progressively replaced by nested attribute-based constraints.
This design enforces structured reasoning over interconnected database records and prevents direct lookup.

\noindent{\textbf{Ambiguous Literature Retrieval.}} In T2, experts apply \emph{Feature Denoising and Fuzzification} to deliberately obscure surface-level cues, while preserving answer uniqueness through \emph{Validation Constraint Injection}, which enables post-retrieval verification without guiding the search process.

\noindent{\textbf{Missing Reference Completion.}} For T3 tasks, we construct citation reconstruction tasks that require aligning scientific claims with their appropriate supporting references.
Surface-level textual overlap is minimized, ensuring that successful prediction depends on understanding claim–evidence relationships rather than keyword matching.

\noindent{\textbf{Cross-Source Structured Knowledge Synthesis.} Tasks in T4 are derived from expert-curated comparison schema and require agents to integrate heterogeneous evidence across multiple papers into a strictly formatted table.
This setting emphasizes cross-source consistency, schema compliance, and high-precision synthesis.}

A detailed description of the construction procedure and validation criteria for each task is provided in Appendix~\ref{app:datasets_construction}.


% \paragraph{Multi-hop Database Navigation.}
% This task evaluates structured reasoning over interconnected scientific databases. Given a query that specifies a target entity indirectly through multiple attribute-based constraints, the agent must traverse several database entries to recover the unique target. Successful completion requires stepwise logical navigation rather than direct keyword lookup.

% \paragraph{Ambiguous Literature Retrieval.}
% This task models exploratory literature search under vague or noisy descriptions. The agent is given an imprecise methodological or conceptual description corresponding to a single target paper and must identify the correct work through broad search and post-retrieval verification. The output is a unique bibliographic record.

% \paragraph{Contextual Citation Reconstruction.}
% This task assesses an agent’s ability to infer supporting literature from scientific context. Given a paragraph with missing citations, the agent must analyze the claims and retrieve the appropriate references that substantiate each claim, relying on contextual reasoning rather than surface-level matching.

% \paragraph{Comparative Knowledge Synthesis.}
% This task represents domain-level synthesis. Given a research topic and a predefined comparison schema, the agent must identify relevant papers, extract heterogeneous methodological details, and synthesize them into a structured comparison table. Evaluation emphasizes both entity coverage and the accuracy of cross-source integration.

% \noindent
% Together, these tasks form a progressive evaluation framework, spanning from entity-level database navigation to document identification, contextual evidence reasoning, and domain-wide structured knowledge synthesis.


% \subsection{Dataset Construction}

% The construction of \textsc{SciExplore} follows a principled, expert-driven methodology designed to faithfully reflect the cognitive demands faced by autonomous scientific research assistants. All task instances are manually crafted by domain experts to ensure realism, difficulty, and alignment with authentic scientific workflows. Rather than relying on automatic generation, our construction process emphasizes controllability, answer uniqueness, and resistance to shortcut retrieval.

% \noindent \textbf{Multi-hop Database Navigation.}
% To construct tasks requiring multi-step structured reasoning, we adopt a \emph{Reverse Trajectory Construction Strategy}. Task generation begins from a uniquely identifiable target entity and progressively replaces explicit entity mentions with composite attribute-based descriptions extracted from scientific database records. During this process, related entities are recursively masked and substituted by their own attributes, forming nested constraints that encode deep dependency chains. Consequently, the final query can only be resolved by reversing the construction trajectory through stepwise database navigation, enforcing genuine multi-hop reasoning and preventing direct lookup.

% \noindent \textbf{Ambiguous Literature Retrieval.}
% Tasks involving ambiguous literature search are constructed to simulate early-stage exploratory research under uncertainty. Experts first extract core methodological or conceptual signals from a target paper and deliberately degrade them into vague, noisy descriptions through \emph{Feature Denoising and Fuzzification}, substantially enlarging the candidate retrieval space. To ensure answer uniqueness despite this ambiguity, we introduce \emph{Validation Constraint Injection}, adding auxiliary constraints that cannot directly guide retrieval but serve as strong post-retrieval validators (e.g., publication venue or author-level properties). This design compels agents to balance broad exploration with rigorous verification.

% \noindent \textbf{Contextual Citation Reconstruction.}
% For tasks requiring citation inference, we select citation-dense paragraphs from high-impact scientific literature and systematically rewrite the surrounding context to remove surface-level textual overlap with the original references. Overly canonical or survey-style citations are excluded to avoid trivial matching. This construction emphasizes reasoning over claim–evidence relationships, ensuring that correct reference reconstruction depends on contextual understanding rather than keyword similarity.

% \noindent \textbf{Comparative Knowledge Synthesis.}
% Tasks targeting structured synthesis are derived from expert-curated comparison tables in review papers across natural sciences and artificial intelligence. Given a predefined comparison schema, experts identify representative source papers and design tasks that require extracting heterogeneous methodological details scattered across multiple documents. Particular care is taken to include variations in terminology, implicit descriptions, and inconsistent reporting conventions. Strict schema and format constraints are enforced so that successful completion depends on accurate cross-source integration rather than partial or approximate extraction.

\begin{table*}[t]
\centering
\caption{Main results on the \textsc{SciExplore}. T4 is reported at two granularities: item- and row-level. All scores are percentages (\%).}
% \vspace{-0.5em}
\small
\label{tabs:main_results}

\begin{tabular*}{\textwidth}{l @{\extracolsep{\fill}} ccccc cc c}
\toprule

% \multirow{2}{*}{\textbf{Model}} & \multirow{2}{*}{\textbf{Category}} &
% \multirow{2}{*}{\textbf{T1}} &
% \multirow{2}{*}{\textbf{T2}} & \multirow{2}{*}{\textbf{T3}} &
% \multicolumn{2}{c}{\textbf{T4}} &
% \multirow{2}{*}{\textbf{Overall}} \\
\multirow{2}{*}[-0.8ex]{\textbf{Model}} & \multirow{2}{*}[-0.8ex]{\textbf{Category}} &
\multirow{2}{*}[-0.8ex]{\textbf{T1}} &
\multirow{2}{*}[-0.8ex]{\textbf{T2}} & \multirow{2}{*}[-0.8ex]{\textbf{T3}} &
\multicolumn{2}{c}{\textbf{T4}} &
\multirow{2}{*}[-0.8ex]{\textbf{Overall}} \\

\cmidrule(lr){6-7}
& & & & & {\scriptsize \textbf{Item}} & {\scriptsize \textbf{Row}} & \\
\midrule

\rowcolor{gray!15}
\multicolumn{8}{l}{\textit{Foundation LLMs}} \\

DeepSeek-V3.2
& Open-Source
& 20.51 & 6.25 & 5.95
& 8.32 & 2.25
& 11.06 \\

Qwen2.5-72B-Instruct
& Open-Source
& 7.69 & 0.00 & 0.00
& 8.63 & 6.94
& 4.99 \\

Qwen3-30B-A3B-Thinking
& Open-Source
& 10.26 & 0.00 & 3.57
& 2.54 & 1.60
& 4.50 \\

Qwen3-235B-A22B-Thinking
& Open-Source
& 20.51 & 3.12 & 0.00
& 7.85 & 4.15
& 7.87 \\

GPT-5.1
& Closed-Source
& 17.95 & 9.38 & 9.66
& 8.76 & 4.83
& 12.83 \\

Gemini-2.5-Pro
& Closed-Source
& 20.51 & 6.25 & 4.17
& 10.57 & 3.78
& 11.25 \\

Gemini-3-Pro
& Closed-Source
& 23.08 & 12.50 & 13.69
& 16.23 & 7.29
& 19.08 \\

\midrule
\rowcolor{gray!15}
\multicolumn{8}{l}{\textit{LLMs with Search Tools}} \\

DeepSeek-V3.2 w/ search
& Open-Source
& 17.95 & 9.38 & 9.66
& 9.20 & 3.95
& 11.12 \\

GPT-5.1 w/ search
& Closed-Source
& 15.38 & 34.38 & 20.85
& 12.13 & 5.12
& 21.61 \\

Gemini-3-Pro w/ search
& Closed-Source
& 33.33 & 53.12 & 49.40
& 23.53 & 11.49
& 46.46 \\

\midrule
\rowcolor{gray!15}
\multicolumn{8}{l}{\textit{Deep Research Agents}} \\

Tongyi-DeepResearch-30B-A3B
& Open-Source
& \textbf{43.59} & 50.00 & 34.40
& 11.12 & 4.26
& 36.93 \\

Gemini Deep Research
& Closed-Source
& 35.90 & \textbf{56.25} & 36.56
& 15.61 & 7.35
& 39.30 \\

OpenAI Deep Research
& Closed-Source
& 23.08 & 43.75 & \textbf{59.91}
& \textbf{30.14} & \textbf{18.59}
& \textbf{49.39} \\

\bottomrule
\end{tabular*}
% \vspace{-1em}
\end{table*}


\subsection{Quality Control}\label{sec:quality_control}

To ensure the rigor of \textsc{SciExplore}, we implemented a streamlined three-stage quality control process: query difficulty calibration, human cross-validation and answer uniqueness verification.

\noindent{\textbf{Query Difficulty Calibration.}}
To minimize parametric memory reliance, we strictly enforce two principles: \textit{Explicit Index Blocking} (removing unique identifiers to compel semantic matching) and \textit{Long-tail Entity Preference} (prioritizing obscure entities). This ensures tasks require deep retrieval reasoning rather than simple knowledge recall.

\noindent{\textbf{Human Cross-Validation.}} 
Annotators independently attempt to solve each question using search engines.
Questions are retained only if no solution can be found within a strict 10-minute limit, ensuring sufficient difficulty for the target models.

\noindent{\textbf{Answer Uniqueness Verification.}} 
We utilize search-enabled LLMs (e.g., Gemini-3-Pro, o3) to generate candidate solutions, which are then reviewed by human annotators. If any alternative answer is satisfied, the task is deemed ambiguous and discarded.

% After the filtering process, we obtain a high-quality set of 103 tasks, with 39, 32, 14, and 18 tasks in T1–T4, respectively.

\begin{comment}
\subsection{Source Statistics}
Unlike benchmarks grounded in a small set of homogeneous web sources, \textsc{SciExplore} spans over ten scientific sub-domains and integrates 16 authoritative databases with a pronounced long-tail distribution. This design forces models to navigate heterogeneous database structures and domain-specific metadata, thereby evaluating their robustness and generalization in authentic, interdisciplinary scientific workflows. Detailed statistics and distribution analyses are provided in Appendix~\ref{appendix_dataset_distribution}.

\subsection{Evaluation Metrics}

We employ task-specific evaluation metrics tailored to the structural characteristics of each task. T1 and T2 are evaluated using Exact Match accuracy based on semantic equivalence of the final answers. T3 is assessed with citation-level precision, measuring the proportion of correctly predicted references for each instance. For T4, we adopt a multi-granularity evaluation protocol at the row and item levels: the row-level metric measures recall over complete table rows, while the item-level metric evaluates fine-grained recall over individual cells or data items. 
An overall score is computed as a weighted sum of the four task scores to reflect holistic model performance. Details are provided in Appendix~\ref{app:evaluation_details}.
\end{comment}